\FigureLayoutDeclare{img/2006/chongqing_liberal/q22_fig1.png}{8.0cm}{0bp 0bp 0bp 0bp}

\chapter{2006年重庆卷（文）}
\section{单选题}
\begin{problem}
已知集合 \(U = \{1, 2, 3, 4, 5, 6, 7\}\)，\(A = \{2, 4, 5, 7\}\)，\(B = \{3, 4, 5\}\)，则 \(\left(\complement_U A\right) \cup \left(\complement_U B\right) =\)（\quad）

\choices
    {\(\{1, 6\}\)}
    {\(\{4, 5\}\)}
    {\(\{2, 3, 4, 5, 7\}\)}
    {\(\{1, 2, 3, 6, 7\}\)}
\end{problem}

\begin{answer}
D
\end{answer}

\begin{solution}
由全集 \(U\) 可得 \(\complement_U A=\{1,3,6\}\)，\(\complement_U B=\{1,2,6,7\}\)，所以
\(\left(\complement_U A\right)\cup\left(\complement_U B\right)=\{1,2,3,6,7\}\)，故选 D.
\end{solution}

\begin{problem}
在等比数列 \(\{a_{n}\}\) 中，若 \(a_{n} > 0\) 且 \(a_{3} a_{7} = 64\)，则 \(a_{5}\) 的值为（\quad）

\choices
    {\(2\)}
    {\(4\)}
    {\(6\)}
    {\(8\)}
\end{problem}

\begin{answer}
D
\end{answer}

\begin{solution}
设等比数列的公比为 \(q\)，则 \(a_3=a_5q^{-2}\)，\(a_7=a_5q^2\)，从而 \(a_3a_7=a_5^2=64\). 又因为 \(a_n>0\)，所以 \(a_5>0\)，故 \(a_5=8\)，选 D.
\end{solution}

\begin{problem}
以点 \((2, -1)\) 为圆心且与直线 \(3x - 4y + 5 = 0\) 相切的圆的方程为（\quad）

\choices
    {\((x - 2)^{2} + (y + 1)^{2} = 3\)}
    {\((x + 2)^{2} + (y - 1)^{2} = 3\)}
    {\((x - 2)^{2} + (y + 1)^{2} = 9\)}
    {\((x + 2)^{2} + (y - 1)^{2} = 9\)}
\end{problem}

\begin{answer}
C
\end{answer}

\begin{solution}
圆心 \((2,-1)\) 到直线 \(3x-4y+5=0\) 的距离为
\(\dfrac{|3\times2-4\times(-1)+5|}{\sqrt{3^2+(-4)^2}}=\dfrac{15}{5}=3\)，所以圆的半径为 \(3\). 因此圆的方程为 \((x-2)^2+(y+1)^2=9\)，选 C.
\end{solution}

\begin{problem}
若 \(P\) 是平面 \(\alpha\) 外一点，则下列命题正确的是（\quad）

\choices
    {过 \(P\) 只能作一条直线与平面 \(\alpha\) 相交}
    {过 \(P\) 可作无数条直线与平面 \(\alpha\) 垂直}
    {过 \(P\) 只能作一条直线与平面 \(\alpha\) 平行}
    {过 \(P\) 可作无数条直线与平面 \(\alpha\) 平行}
\end{problem}

\begin{answer}
D
\end{answer}

\begin{solution}
过平面 \(\alpha\) 外一点 \(P\) 的垂线方向由平面的法向量唯一确定，所以只能作一条直线垂直于 \(\alpha\). 过 \(P\) 的直线若与 \(\alpha\) 相交，可以改变相交点，因而不只有一条；平行于 \(\alpha\) 的直线方向可以取平面内任意方向，故有无数条. 所以正确命题是 D.
\end{solution}

\begin{problem}
\((2x - 3)^5\) 的展开式中 \(x^{2}\) 项的系数为（\quad）

\choices
    {\(-2160\)}
    {\(-1080\)}
    {\(1080\)}
    {\(2160\)}
\end{problem}

\begin{answer}
B
\end{answer}

\begin{solution}
二项式通项中含 \(x^2\) 的项必须从五个因式中取两个 \(2x\)，其余三个取 \(-3\)，所以该项为 \(\mathrm C_5^2(2x)^2(-3)^3\). 因而 \(x^2\) 项的系数为 \(10\times4\times(-27)=-1080\)，选 B.
\end{solution}

\begin{problem}
设函数 \( y = f(x) \) 的反函数为 \( y = f^{-1}(x) \)，且 \( y = f(2x - 1) \) 的图象过点 \( \left(\frac{1}{2}, 1\right) \)，则 \( y = f^{-1}(x) \) 的图象必过点（\quad）

\choices
    {\(\left(\frac{1}{2}, 1\right)\)}
    {\(\left(1, \frac{1}{2}\right)\)}
    {\((1,0)\)}
    {\((0,1)\)}
\end{problem}

\begin{answer}
C
\end{answer}

\begin{solution}
由 \(y=f(2x-1)\) 的图象经过 \(\left(\frac12,1\right)\)，代入得 \(f(0)=1\). 因为反函数图象上的点由原函数图象上的点交换横、纵坐标得到，所以 \(y=f^{-1}(x)\) 的图象必经过 \((1,0)\)，选 C.
\end{solution}

\begin{problem}
某地区有 \(300\) 家商店，其中大型商店有 \(30\) 家，中型商店有 \(75\) 家，小型商店有 \(195\) 家. 为了掌握各商店的营业情况，要从中抽取一个容量为 \(20\) 的样本. 若采用分层抽样的方法，抽取的中型商店数是（\quad）

\choices
    {\(2\)}
    {\(3\)}
    {\(5\)}
    {\(13\)}
\end{problem}

\begin{answer}
C
\end{answer}

\begin{solution}
中型商店占全部商店的比例为 \(\dfrac{75}{300}\). 按分层抽样，容量为 \(20\) 的样本中应抽取中型商店
\(20\times\dfrac{75}{300}=5\) 家，选 C.
\end{solution}

\begin{problem}
已知三点 \(A(2,3)\)，\(B(-1,-1)\)，\(C(6,k)\). 其中 \(k\) 为常数. 若 \(\left|AB\right| = \left|AC\right|\)，则 \(\overrightarrow{AB}\) 与 \(\overrightarrow{AC}\) 的夹角为（\quad）

\choices
    {\(\arccos \left(-\frac{24}{25}\right)\)}
    {\(\frac{\pi}{2}\) 或 \(\arccos \frac{24}{25}\)}
    {\(\arccos \frac{24}{25}\)}
    {\(\frac{\pi}{2}\) 或 \(\pi -\arccos \frac{24}{25}\)}
\end{problem}

\begin{answer}
D
\end{answer}

\begin{solution}
\(\overrightarrow{AB}=(-3,-4)\)，\(|AB|=5\)，而 \(\overrightarrow{AC}=(4,k-3)\). 由 \(|AB|=|AC|\) 得
\(16+(k-3)^2=25\)，所以 \(k=6\) 或 \(k=0\).

当 \(k=6\) 时，\(\overrightarrow{AC}=(4,3)\)，于是
\(\overrightarrow{AB}\cdot\overrightarrow{AC}=-3\times4-4\times3=-24\)，夹角 \(\theta\) 满足 \(\cos\theta=-\dfrac{24}{25}\)，即 \(\theta=\pi-\arccos\dfrac{24}{25}\). 当 \(k=0\) 时，\(\overrightarrow{AC}=(4,-3)\)，有 \(\overrightarrow{AB}\cdot\overrightarrow{AC}=0\)，故 \(\theta=\dfrac\pi2\). 所以选 D.
\end{solution}

\begin{problem}
高三（一）班需要安排毕业晚会的\(4\)个音乐节目，\(2\)个舞蹈节目和\(1\)个曲艺节目的演出顺序，要求两个舞蹈节目不连排，则不同排法的种数是（\quad）

\choices
    {\(1800\)}
    {\(3600\)}
    {\(4320\)}
    {\(5040\)}
\end{problem}

\begin{answer}
B
\end{answer}

\begin{solution}
先把两个舞蹈节目看作一个整体，则共有 \(6!\) 种排法；在这个整体内部，两个舞蹈节目的顺序有 \(2!\) 种，所以相连的排法有 \(2\times6!=1440\) 种. 全部七个不同节目的排法有 \(7!=5040\) 种，故两个舞蹈节目不连排的排法有
\(7!-2\times6!=5040-1440=3600\) 种，选 B.
\end{solution}

\begin{problem}
若 \(\alpha\)，\(\beta \in \left(0, \frac{\pi}{2}\right)\)，\(\cos \left(\alpha - \frac{\beta}{2}\right) = \frac{\sqrt{3}}{2}\)，\(\sin \left(\frac{\alpha}{2} - \beta\right) = -\frac{1}{2}\)，则 \(\cos (\alpha + \beta)\) 的值等于（\quad）

\choices
    {\(-\frac{\sqrt{3}}{2}\)}
    {\(-\frac{1}{2}\)}
    {\( \frac{1}{2} \)}
    {\(\frac{\sqrt{3}}{2}\)}
\end{problem}

\begin{answer}
B
\end{answer}

\begin{solution}
记 \(u=\alpha-\dfrac\beta2\)，\(v=\dfrac\alpha2-\beta\). 由 \(0<\alpha\)，\(\beta<\dfrac\pi2\)，知 \(-\dfrac\pi4<u<\dfrac\pi2\)，\(-\dfrac\pi2<v<\dfrac\pi4\). 因而由
\(\cos u=\dfrac{\sqrt3}{2}\) 得 \(u=\dfrac\pi6\) 或 \(u=-\dfrac\pi6\)，由 \(\sin v=-\dfrac12\) 得 \(v=-\dfrac\pi6\).

若 \(u=\dfrac\pi6\)，则由 \(\alpha-\dfrac\beta2=\dfrac\pi6\) 与 \(\dfrac\alpha2-\beta=-\dfrac\pi6\) 解得 \(\alpha=\beta=\dfrac\pi3\). 若 \(u=-\dfrac\pi6\)，联立两式解得 \(\alpha=-\dfrac\pi9\)，不满足 \(\alpha>0\)，舍去. 因此 \(\cos(\alpha+\beta)=\cos\dfrac{2\pi}{3}=-\dfrac12\)，选 B.
\end{solution}

\begin{problem}
设 \(A(x_{1},y_{1})\)，\(B\left(4, \frac{9}{5}\right)\)，\(C(x_{2},y_{2})\) 是右焦点为 \(F\) 的椭圆 \(\frac{x^2}{25} + \frac{y^2}{9} = 1\) 上三个不同的点，则 “\(|AF|\)，\(|BF|\)，\(|CF|\) 成等差数列”是 “\(x_{1} + x_{2} = 8\)” 的（\quad）

\choices
    {充要条件}
    {必要而不充分条件}
    {充分而不必要条件}
    {既不充分也不必要条件}
\end{problem}

\begin{answer}
A
\end{answer}

\begin{solution}
设椭圆的左焦点为 \(F'(-4,0)\)，右焦点为 \(F(4,0)\). 对椭圆上任一点 \(P(x,y)\)，由椭圆定义有 \(|PF|+|PF'|=10\). 又
\(|PF'|^2-|PF|^2=((x+4)^2+y^2)-((x-4)^2+y^2)=16x\)，所以
\(|PF'|-|PF|=\dfrac{8x}{5}\)，从而 \(|PF|=5-\dfrac{4x}{5}\).

对点 \(B\)，\(|BF|=\dfrac95\). 因此
\(|AF|\)，\(|BF|\)，\(|CF|\) 成等差数列，当且仅当 \(2|BF|=|AF|+|CF|\)，即
\(\dfrac{18}{5}=10-\dfrac45(x_1+x_2)\)，这等价于 \(x_1+x_2=8\). 故二者互为充要条件，选 A.
\end{solution}

\begin{problem}
若 \(a\)，\(b\)，\(c > 0\) 且 \(a^2 + 2ab + 2ac + 4bc = 12\)，则 \(a + b + c\) 的最小值是（\quad）

\choices
    {\(2\sqrt{3}\)}
    {\(3\)}
    {\(2\)}
    {\( \sqrt{3} \)}
\end{problem}

\begin{answer}
A
\end{answer}

\begin{solution}
记 \(S=a+b+c\). 由
\(S^2=a^2+b^2+c^2+2ab+2ac+2bc\)，可得
\(a^2+2ab+2ac+4bc=S^2-(b-c)^2\). 题设于是给出
\(S^2=12+(b-c)^2\geqslant12\). 因为 \(S>0\)，所以 \(S\geqslant2\sqrt3\).

当 \(b=c=\dfrac{\sqrt3}{2}\)，\(a=\sqrt3\) 时，题设等式成立且 \(a+b+c=2\sqrt3\)，故最小值为 \(2\sqrt3\)，选 A.
\end{solution}

\section{填空题}
\begin{problem}
已知 \(\sin \alpha = \frac{2\sqrt{5}}{5}\)，\(\frac{\pi}{2} < \alpha < \pi\)，则 \(\tan \alpha = \)\fillinblank{}.
\end{problem}

\begin{answer}
\(-2\)
\end{answer}

\begin{solution}
由 \(\dfrac\pi2<\alpha<\pi\)，知 \(\cos\alpha<0\). 由同角三角函数基本关系，
\(\cos\alpha=-\sqrt{1-\sin^2\alpha}=-\sqrt{1-\dfrac45}=-\dfrac{\sqrt5}{5}\). 因而
\(\tan\alpha=\dfrac{\sin\alpha}{\cos\alpha}=\dfrac{2\sqrt5/5}{-\sqrt5/5}=-2\).
\end{solution}

\begin{problem}
在数列 \( \{a_{n}\} \) 中，若 \(a_{1}=1\)，\(a_{n+1}=a_{n}+2\)（\(n\geqslant1\)），则该数列的通项 \( a_{n}= \)\fillinblank{}.
\end{problem}

\begin{answer}
\(2n-1\)
\end{answer}

\begin{solution}
由递推式 \(a_{n+1}-a_n=2\)，数列 \(\{a_n\}\) 是首项为 \(1\)、公差为 \(2\) 的等差数列. 所以
\(a_n=a_1+(n-1)\times2=1+2(n-1)=2n-1\).
\end{solution}

\begin{problem}
设 \(a > 0\)，\(a \neq 1\)，函数 \(f(x) = \log_a(x^2 - 2x + 3)\) 有最小值，则不等式 \(\log_a(x - 1) > 0\) 的解集为\fillinblank{}.
\end{problem}

\begin{answer}
\((2,+\infty)\)
\end{answer}

\begin{solution}
函数的真数为 \(x^2-2x+3=(x-1)^2+2\)，其最小值为 \(2\). 若 \(0<a<1\)，则对数函数在正数上单调递减，且当 \(|x|\to+\infty\) 时 \(\log_a((x-1)^2+2)\to-\infty\)，函数不可能有最小值. 因此题设的“有最小值”说明 \(a>1\).

此时 \(\log_a(x-1)>0\) 要求真数 \(x-1>0\)，并且等价于 \(x-1>1\)，所以解集为 \((2,+\infty)\).
\end{solution}

\begin{problem}
已知变量 \(x\)，\(y\) 满足约束条件 \(\begin{cases}
x + 2y - 3 \leqslant 0,\\
x + 3y - 3 \geqslant 0,\\
y - 1 \leqslant 0,
\end{cases}\) 且若目标函数 \(z = ax + y\)（其中 \(a > 0\)）仅在点 \((3,0)\) 处取得最大值，则 \(a\) 的取值范围为\fillinblank{}.
\end{problem}

\begin{answer}
\(a>\frac{1}{2}\)
\end{answer}

\begin{solution}
由约束条件可写成 \(x\leqslant3-2y\)、\(x\geqslant3-3y\)、\(y\leqslant1\). 要使前两个不等式有公共的 \(x\)，必须有 \(3-3y\leqslant3-2y\)，即 \(y\geqslant0\). 所以可行域是顶点为 \((3,0)\)、\((1,1)\)、\((0,1)\) 的三角形.

目标函数 \(z=ax+y\) 在三个顶点处的值分别为 \(3a\)、\(a+1\)、\(1\). 要使它仅在 \((3,0)\) 处取得最大值，必须有 \(3a>a+1\) 且 \(3a>1\). 因为 \(a>0\)，前一个条件给出 \(a>\dfrac12\)，并且它已蕴含后一个条件. 反之，当 \(a>\dfrac12\) 时，\((3,0)\) 的函数值严格大于另外两个顶点的函数值，故最大值仅在该点取得. 因此 \(a\) 的取值范围为 \(a>\dfrac12\).
\end{solution}

\section{解答题}
\begin{problem}
甲，乙，丙三人在同一办公室工作，办公室里只有一部电话机. 设该机打进的电话是打给甲，乙，丙的概率依次为 \(\frac{1}{6}\)，\(\frac{1}{3}\)，\(\frac{1}{2}\). 若在一段时间内打进三个电话，且各个电话相互独立. 求：

\begin{enumerate}
\item 这三个电话是打给同一个人的概率；
\item 这三个电话中恰有两个是打给甲的概率.
\end{enumerate}
\end{problem}

\begin{solution}
设三个电话分别记为第一个、第二个和第三个电话.

\begin{enumerate}
\item 三个电话都打给甲、乙、丙的概率分别为 \(\left(\frac{1}{6}\right)^3\)、\(\left(\frac{1}{3}\right)^3\)、\(\left(\frac{1}{2}\right)^3\)，所以三个电话打给同一个人的概率为 \(\left(\frac{1}{6}\right)^3+\left(\frac{1}{3}\right)^3+\left(\frac{1}{2}\right)^3=\frac{1}{216}+\frac{8}{216}+\frac{27}{216}=\frac{1}{6}\).
\item 恰有两个电话打给甲时，先从三个电话中选出两个打给甲，另一个电话不能打给甲，故所求概率为 \(\mathrm C_3^2\left(\frac{1}{6}\right)^2\left(1-\frac{1}{6}\right)=3\cdot\frac{1}{36}\cdot\frac{5}{6}=\frac{5}{72}\).
\end{enumerate}
\end{solution}

\begin{problem}
设函数 \( f(x) = \sqrt{3}\cos^2\omega x + \sin \omega x\cos \omega x + a \) （其中 \(\omega > 0\)，\(a \in \R\)），且 \( f(x) \) 的图象在 \( y \) 轴右侧的第一个最高点的横坐标为 \( \frac{\pi}{6} \).

\begin{enumerate}
\item 求 \(\omega\) 的值；
\item 如果 \( f(x) \) 在区间 \(\left[-\frac{\pi}{3}, \frac{5\pi}{6}\right]\) 上的最小值为 \(\sqrt{3}\)，求 \(a\) 的值.
\end{enumerate}
\end{problem}

\begin{solution}
利用二倍角公式，有 \(\sqrt{3}\cos^2\omega x=\frac{\sqrt{3}}{2}+\frac{\sqrt{3}}{2}\cos 2\omega x\)，\(\sin\omega x\cos\omega x=\frac{1}{2}\sin 2\omega x\)，因此
\[
f(x)=a+\frac{\sqrt{3}}{2}+\frac{\sqrt{3}}{2}\cos 2\omega x+\frac{1}{2}\sin 2\omega x=a+\frac{\sqrt{3}}{2}+\cos\left(2\omega x-\frac{\pi}{6}\right)
\]

\begin{enumerate}
\item 当 \(2\omega x-\frac{\pi}{6}=2k\pi\) 时，\(f(x)\) 取得最高值. 因为 \(\omega>0\)，在 \(y\) 轴右侧的第一个最高点对应于 \(k=0\)，其横坐标为 \(\frac{\pi}{12\omega}\). 由题意 \(\frac{\pi}{12\omega}=\frac{\pi}{6}\)，解得 \(\omega=\frac{1}{2}\).
\item 由 \(\omega=\frac{1}{2}\)，得 \(f(x)=a+\frac{\sqrt{3}}{2}+\cos\left(x-\frac{\pi}{6}\right)\). 当 \(x\in\left[-\frac{\pi}{3},\frac{5\pi}{6}\right]\) 时，\(x-\frac{\pi}{6}\in\left[-\frac{\pi}{2},\frac{2\pi}{3}\right]\). 在这个区间上，余弦函数的最小值为 \(\cos\frac{2\pi}{3}=-\frac{1}{2}\)，所以 \(f(x)\) 的最小值为 \(a+\frac{\sqrt{3}}{2}-\frac{1}{2}\). 由题意得 \(a+\frac{\sqrt{3}}{2}-\frac{1}{2}=\sqrt{3}\)，故 \(a=\frac{\sqrt{3}+1}{2}\).
\end{enumerate}
\end{solution}

\begin{problem}
设函数 \( f(x) = x^3 - 3ax^2 + 3bx \) 的图象与直线 \( 12x + y - 1 = 0 \) 相切于点 \((1,-11)\).

\begin{enumerate}
\item 求 \(a\)，\(b\) 的值；
\item 讨论函数 \( f(x) \) 的单调性.
\end{enumerate}
\end{problem}

\begin{solution}
\begin{enumerate}
\item 因为切点 \(\left(1,-11\right)\) 在函数图象上，所以 \(f(1)=-11\)，即 \(1-3a+3b=-11\). 又因为直线 \(12x+y-1=0\) 的斜率为 \(-12\)，且该直线与函数图象在 \(\left(1,-11\right)\) 处相切，所以 \(f'(1)=-12\). 由 \(f'(x)=3x^2-6ax+3b\)，得
\[
\begin{cases}
1-3a+3b=-11,\\
3-6a+3b=-12.
\end{cases}
\]
两式分别化简为 \(b-a=-4\) 和 \(b-2a=-5\)，解得 \(a=1\)，\(b=-3\).
\item 此时 \(f'(x)=3x^2-6x-9=3(x+1)(x-3)\). 因而当 \(x\in(-\infty,-1)\) 或 \(x\in(3,+\infty)\) 时，\(f'(x)>0\)，函数 \(f(x)\) 单调递增；当 \(x\in(-1,3)\) 时，\(f'(x)<0\)，函数 \(f(x)\) 单调递减.
\end{enumerate}
\end{solution}

\begin{problem}
如图，在正四棱柱 \(ABCD - A_{1}B_{1}C_{1}D_{1}\) 中，\(AB = 1\)，\(BB_{1} = \sqrt{3} + 1\)，\(E\) 为 \(BB_{1}\) 上使 \(B_{1}E = 1\) 的点，平面 \(AEC_{1}\) 交 \(DD_{1}\) 于 \(F\)，交 \(A_{1}D_{1}\) 的延长线于 \(G\). 求：

\begin{enumerate}
\item 异面直线 \(AD\) 与 \(C_{1}G\) 所成角的大小；
\item 二面角 \(A - C_{1}G - A_{1}\) 的正切值.
\end{enumerate}

\bitmapfigure[width=0.30\linewidth]{img/2006/chongqing_liberal/q20_fig1.png}
\end{problem}

\begin{solution}
取 \(B_{1}\) 为原点，令 \(B_{1}C_{1}\)、\(B_{1}A_{1}\)、\(B_{1}B\) 分别为三个坐标轴的正方向，记 \(h=\sqrt{3}+1\). 则 \(A_{1}(0,1,0)\)，\(C_{1}(1,0,0)\)，\(D_{1}(1,1,0)\)，\(A(0,1,h)\)，\(D(1,1,h)\)，\(E(0,0,1)\).

由 \(\overrightarrow{EC_{1}}=(1,0,-1)\)、\(\overrightarrow{EA}=(0,1,\sqrt{3})\) 可得平面 \(AEC_{1}\) 的一个法向量为 \((1,-\sqrt{3},1)\)，所以其方程为 \(x-\sqrt{3}y+z-1=0\). 在直线 \(DD_{1}\) 上有 \(x=y=1\)，代入平面方程得 \(z=\sqrt{3}\)，所以 \(F=(1,1,\sqrt{3})\). 在直线 \(A_{1}D_{1}\) 的延长线上有 \(y=1\)、\(z=0\)，代入平面方程得 \(x=\sqrt{3}+1=h\)，所以 \(G=(h,1,0)\).

\begin{enumerate}
\item \(\overrightarrow{AD}=(1,0,0)\)，\(\overrightarrow{C_{1}G}=(h-1,1,0)=(\sqrt{3},1,0)\). 设异面直线 \(AD\) 与 \(C_{1}G\) 所成的角为 \(\theta\)，则 \(\cos\theta=\frac{\sqrt{3}}{\sqrt{3+1}}=\frac{\sqrt{3}}{2}\). 因而 \(\theta=\frac{\pi}{6}\).
\item 平面 \(AC_{1}G\) 即平面 \(AEC_{1}\)，其一个法向量为 \((1,-\sqrt{3},1)\)；平面 \(A_{1}C_{1}G\) 是底面，其一个法向量为 \((0,0,1)\). 设二面角的平面角为 \(\varphi\)，则两个法向量的夹角满足 \(\cos\varphi=\frac{1}{\sqrt{5}}\)，从而 \(\tan\varphi=\frac{\sqrt{1-\frac{1}{5}}}{\frac{1}{\sqrt{5}}}=2\). 因此二面角 \(A-C_{1}G-A_{1}\) 的正切值为 \(2\).
\end{enumerate}
\end{solution}

\begin{problem}
已知定义域为 \(\R\) 的函数 \(f(x) = \frac{-2^x + b}{2^{x + 1} + a}\) 是奇函数.

\begin{enumerate}
\item 求 \(a\)，\(b\) 的值；
\item 若对任意的 \(t \in \R\)，不等式 \(f(t^2 - 2t) + f(2t^2 - k) < 0\) 恒成立，求 \(k\) 的取值范围.
\end{enumerate}
\end{problem}

\begin{solution}
\begin{enumerate}
\item 令 \(u=2^x\)，则 \(u>0\)，且 \(f(x)=\frac{b-u}{2u+a}\). 由 \(f(x)\) 的定义域为 \(\R\)，知 \(2u+a\) 对一切 \(u>0\) 均不为零. 又因为 \(f(x)\) 是奇函数，所以
\[
\frac{bu-1}{2+au}=f(-x)=-f(x)=\frac{u-b}{2u+a}
\]
于是对一切 \(u>0\)，有 \(\left(bu-1\right)\left(2u+a\right)=\left(u-b\right)\left(2+au\right)\)，即 \(2bu^{2}+(ab-2)u-a=au^{2}+(2-ab)u-2b\). 比较两边关于 \(u\) 的系数，得到 \(a=2b\) 且 \(ab=2\)，所以 \(b=\pm1\)，\(a=\pm2\)，且 \(a\)，\(b\) 同号. 当 \(a=-2\) 时，\(x=0\) 使分母 \(2^{x+1}+a=0\)，与定义域为 \(\R\) 矛盾，故 \(a=2\)，\(b=1\).
\item 代入 \(a=2\)，\(b=1\)，得 \(f(x)=\frac{1-2^x}{2(1+2^x)}\). 令 \(u=2^{t^2-2t}\)，\(v=2^{2t^2-k}\)，则 \(u>0\)，\(v>0\)，并且
\[
f(t^2-2t)+f(2t^2-k)=\frac{1-uv}{(1+u)(1+v)}
\]
分母恒为正，因此原不等式等价于 \(uv>1\)，即 \(2^{3t^2-2t-k}>1\)，也即 \(3t^2-2t-k>0\) 对任意 \(t\in\R\) 恒成立. 配方得 \(3t^2-2t-k=3\left(t-\frac{1}{3}\right)^2-\frac{1}{3}-k\)，其最小值为 \(-\frac{1}{3}-k\). 由于题设要求严格大于零，必须且只需 \(-\frac{1}{3}-k>0\)，故 \(k<-\frac{1}{3}\). 当 \(k=-\frac{1}{3}\) 时在 \(t=\frac{1}{3}\) 处等号成立，不能取；当 \(k>-\frac{1}{3}\) 时在 \(t=\frac{1}{3}\) 处不等式更不成立.
\end{enumerate}
\end{solution}

\begin{problem}
如图，对每个正整数 \(n\)，\(A_{n}(x_{n}, y_{n})\) 是抛物线 \(x^{2} = 4y\) 上的点，过焦点 \(F\) 的直线 \(FA_{n}\) 交抛物线于另一点 \(B_{n}(s_{n}, t_{n})\).

\begin{enumerate}
\item 试证：\(x_{n}s_{n} = -4\)（\(n \geqslant 1\)）；
\item 取 \(x_{n} = 2^{n}\)，并记 \(C_{n}\) 为抛物线上分别以 \(A_{n}\) 与 \(B_{n}\) 为切点的两条切线的交点. 试证：\(|FC_1| + |FC_2| + \cdots + |FC_n| = 2^n - 2^{-n+1} + 1\).
\end{enumerate}

\bitmapfigure[width=0.42\linewidth]{img/2006/chongqing_liberal/q22_fig1.png}
\end{problem}

\begin{solution}
抛物线 \(x^2=4y\) 的焦点为 \(F(0,1)\).

\begin{enumerate}
\item 由于直线 \(FA_n\) 与抛物线有两个不同交点，它不可能是竖直直线，故可设为 \(y-1=mx\). 将它与抛物线方程联立，得 \(x^2=4(mx+1)\)，即 \(x^2-4mx-4=0\). 这个方程的两个根正是 \(x_n\) 和 \(s_n\)，由韦达定理得 \(x_ns_n=-4\).
\item 记 \(u=x_n\)，\(v=s_n\). 抛物线在横坐标为 \(r\) 的点 \(\left(r,\frac{r^2}{4}\right)\) 处的切线方程为 \(y-\frac{r^2}{4}=\frac{r}{2}(x-r)\)，即 \(y=\frac{r}{2}x-\frac{r^2}{4}\). 联立横坐标分别为 \(u\)、\(v\) 的两条切线，得交点 \(C_n\left(\frac{u+v}{2},\frac{uv}{4}\right)=\left(\frac{u+v}{2},-1\right)\)，其中使用了 \(uv=-4\). 由于 \((u+v)^2+16=u^2+2uv+v^2+16=u^2-2uv+v^2=(u-v)^2\)，所以
\[
|FC_n|=\sqrt{\left(\frac{u+v}{2}\right)^2+(-2)^2}=\frac{|u-v|}{2}
\]
当 \(x_n=2^n\) 时，由第一问 \(s_n=-\frac{4}{2^n}=-2^{2-n}\)，所以 \(u>0>v\)，从而
\[
|FC_n|=\frac{1}{2}\left(2^n+2^{2-n}\right)=2^{n-1}+2^{1-n}
\]
于是
\[
\sum_{j=1}^{n}|FC_j|=\sum_{j=1}^{n}2^{j-1}+\sum_{j=1}^{n}2^{1-j}=(2^n-1)+(2-2^{1-n})=2^n-2^{-n+1}+1
\]
这就证明了所要证的等式.
\end{enumerate}
\end{solution}
